\noindent These are descriptive accounting comparisons, not a causal model of cabinet formation or of party-size effects. Each pooled observation is one election-year party, equally weighted. Cabinet membership uses the existing translated election-year registry. The 5\% benchmark is descriptive only, not an electoral threshold or theoretical cutoff. Since $q_i=513v_i/V$, vote share and $q_i$ have identical correlations. Ratios are dimensionless; $A_i$ is measured in seats.
\begingroup
\scriptsize
\setlength{\tabcolsep}{3.2pt}
\begin{longtable}{llrrrrr}
\caption{Party size by cabinet participation}\label{tab:report-party-size-cabinet} \\
\toprule
Election & Ever cabinet & n & Mean vote \% & Median vote \% & Mean q & Median q \\
\midrule
\endfirsthead
\multicolumn{7}{l}{\footnotesize Party size by cabinet participation (continued)} \\
\toprule
Election & Ever cabinet & n & Mean vote \% & Median vote \% & Mean q & Median q \\
\midrule
\endhead
\midrule
\multicolumn{7}{r}{\footnotesize Continued on next page} \\
\endfoot
\bottomrule
\endlastfoot
2014 & Yes & 15 & 5.63 & 4.55 & 28.91 & 23.32 \\
2014 & No & 17 & 0.91 & 0.69 & 4.67 & 3.52 \\
2018 & Yes & 11 & 5.05 & 5.32 & 25.92 & 27.28 \\
2018 & No & 24 & 1.85 & 1.40 & 9.49 & 7.20 \\
2022 & Yes & 11 & 5.80 & 6.96 & 29.75 & 35.68 \\
2022 & No & 21 & 1.72 & 0.87 & 8.85 & 4.48 \\
pooled & Yes & 37 & 5.51 & 5.32 & 28.27 & 27.28 \\
pooled & No & 62 & 1.55 & 0.83 & 7.95 & 4.27 \\
\end{longtable}
\endgroup

\begingroup
\scriptsize
\setlength{\tabcolsep}{3.2pt}
\begin{longtable}{lrrrrr}
\caption{Party-size associations}\label{tab:report-party-size-correlations} \\
\toprule
Election & n & Pearson A & Pearson A/q & Spearman A & Spearman A/q \\
\midrule
\endfirsthead
\multicolumn{6}{l}{\footnotesize Party-size associations (continued)} \\
\toprule
Election & n & Pearson A & Pearson A/q & Spearman A & Spearman A/q \\
\midrule
\endhead
\midrule
\multicolumn{6}{r}{\footnotesize Continued on next page} \\
\endfoot
\bottomrule
\endlastfoot
2014 & 32 & 0.455 & 0.506 & 0.475 & 0.624 \\
2018 & 35 & 0.520 & 0.647 & 0.408 & 0.781 \\
2022 & 32 & 0.810 & 0.708 & 0.150 & 0.803 \\
pooled & 99 & 0.634 & 0.609 & 0.383 & 0.774 \\
\end{longtable}
\endgroup

\begingroup
\scriptsize
\setlength{\tabcolsep}{3.2pt}
\begin{longtable}{llrrrr}
\caption{Descriptive national-vote size groups}\label{tab:report-party-size-bins} \\
\toprule
Election & Vote group & n & A mean / median & A/q mean / median & Positive A: n (\%) \\
\midrule
\endfirsthead
\multicolumn{6}{l}{\footnotesize Descriptive national-vote size groups (continued)} \\
\toprule
Election & Vote group & n & A mean / median & A/q mean / median & Positive A: n (\%) \\
\midrule
\endhead
\midrule
\multicolumn{6}{r}{\footnotesize Continued on next page} \\
\endfoot
\bottomrule
\endlastfoot
2014 & <1\% & 14 & -1.04 / -0.95 & -0.536 / -0.622 & 3 (21.43) \\
2014 & 1 to <3\% & 7 & -0.74 / -0.33 & -0.080 / -0.025 & 2 (28.57) \\
2014 & 3 to <5\% & 4 & 0.88 / -0.15 & 0.043 / -0.007 & 1 (25.00) \\
2014 & >=5\% & 7 & 2.32 / 2.64 & 0.068 / 0.069 & 6 (85.71) \\
2018 & <1\% & 11 & -1.37 / -1.15 & -0.754 / -0.924 & 0 (0.00) \\
2018 & 1 to <3\% & 13 & -1.98 / -2.76 & -0.195 / -0.293 & 2 (15.38) \\
2018 & 3 to <5\% & 2 & 3.84 / 3.84 & 0.161 / 0.161 & 2 (100.00) \\
2018 & >=5\% & 9 & 3.68 / 4.34 & 0.128 / 0.142 & 8 (88.89) \\
2022 & <1\% & 12 & -0.56 / -0.46 & -0.790 / -0.933 & 1 (8.33) \\
2022 & 1 to <3\% & 8 & -3.01 / -3.14 & -0.394 / -0.416 & 1 (12.50) \\
2022 & 3 to <5\% & 5 & -3.71 / -3.40 & -0.209 / -0.188 & 0 (0.00) \\
2022 & >=5\% & 7 & 7.06 / 5.12 & 0.125 / 0.137 & 7 (100.00) \\
pooled & <1\% & 37 & -0.98 / -0.78 & -0.683 / -0.844 & 4 (10.81) \\
pooled & 1 to <3\% & 28 & -1.97 / -2.48 & -0.223 / -0.264 & 5 (17.86) \\
pooled & 3 to <5\% & 11 & -0.67 / -0.25 & -0.050 / -0.012 & 3 (27.27) \\
pooled & >=5\% & 23 & 4.30 / 4.26 & 0.109 / 0.126 & 21 (91.30) \\
\end{longtable}
\endgroup

Across party-elections, 21/23 parties with at least 5\% of votes have positive \(A_i\), compared with 12/76 below 5\%. Cabinet participants have larger means and medians in each election, and normalization by \(q_i\) retains a positive size association. This is not a monotonic size gradient: the 2022 intermediate-size group (3 to less than 5\%) is entirely negative, while many tiny parties have small absolute losses.
The large-party exception PT 2014 has 13.92\% of votes, \(A_i=-1.22\) and \(A_i/q_i=-0.017\).
The large-party exception PSL 2018 has 11.64\% of votes, \(A_i=-2.42\) and \(A_i/q_i=-0.041\).
Large parties also remain outside cabinet: the largest nonparticipant in each election is SOLIDARIEDADE 2014 (2.76\% of votes); PT 2018 (10.31\% of votes); PL 2022 (16.61\% of votes).
\par All 23 current cabinet observations have positive \(A_C\); \(B_C>0\) in 6 of them. There are 22 distinct election-year party sets. The nonadjacent 2020.4 and 2021.1 observations share a set, but remain separate chronology rows. The existing adjacent 2021.3/2022.1 observation is unchanged. Positive contributions from members with at least 5\% of votes exceed all negative member contributions in every set; they supply 68.67--99.42\% of gross positives. The smallest such remaining balance is 1.10 seats. Net \(A_C\) ranges from 5.68 to 22.14 seats.
\begin{itemize}
\item 2014 recurring positive members: PR: \(A_i=5.04\), 4 observations; PSD: \(A_i=4.65\), 7 observations; PTB: \(A_i=4.08\), 5 observations; PMDB: \(A_i=3.91\), 8 observations.
\item 2018 recurring positive members: PP: \(A_i=7.32\), 4 observations; PR: \(A_i=6.37\), 11 observations; DEM: \(A_i=4.58\), 12 observations; PRB: \(A_i=4.38\), 12 observations.
\item 2022 recurring positive members: UNIÃO: \(A_i=9.36\), 2 observations; PT: \(A_i=9.07\), 3 observations; MDB: \(A_i=5.12\), 3 observations.
\end{itemize}
The weakest net total belongs to period(s) 2020.2: \(A_C=5.68\), with gross positives 15.75 and negatives -10.07. The three largest members by \(q_i\) supply as little as 17.14\% of gross positive \(A_i\); the explanation concerns a broader group of relatively large members. All distinct-set decompositions and all 23 period links are retained in the machine-readable outputs. The group arithmetic holds fixed the observed party contributions; it is not a simulated seat allocation after removing parties.
